\scanpage{54}
\mspage{53}
\unsure{Contrairement|(contrairement} à l'exemple B) p. 34) sans structure affine. On définit la compatibilité des simplexes comme dans loc. cit. \deleted{Mais \unsure{ce n'est pas} c'est}
D'où \unsure{ens} de multistrates, \unsure{on prendra|où pour des} $\unsure{\mathcal F|\mathscr F}$ \unsure{ensemble|en met} des \deleted{multistrates}\inserted{simplexes} fermé\deleted{\unsure{e|es}}s pour $\leq$, et formés de simplexes deux à deux compatibles. L'axiome A 8 \unsure{trivialmt|trivialement} valable (NB \deleted{\unsure{$L$}}$E$ s'identifie à $L$). Mais la notion de compatibilité des simplexes est plus fine que la compatibilité des figures ensemblistes associées -- puisqu'ici on exige \inserted{de plus} que la figure intersection soit élémentaire, i.e. soit une multistrate commune. On y reviendra par la suite\ldots

Ainsi, dans le cas d'un atelier ensembliste, $\mathcal F$, $\mathcal F$ se réconstitue à partir de \deleted{\unsure{un}} l'ens des figures \unsure{dans|des} $\unsure{\mathcal L|L}$ qu'il décrit, comme sous-atelier de $\operatorname{Fig}(L)$. La question est donc \uline{quels} sous-ateliers d'un ens $L$ on obtient ainsi.
